%=============================================================================
% AGENT 10: UAP FLIGHT PROFILE MATCHING AGAINST DFD PSI-BUBBLE CONSTRAINTS
% Systematic comparison of reported UAP capabilities vs DFD physics
%=============================================================================

\documentclass[12pt,a4paper]{article}
\usepackage{amsmath,amssymb,amsfonts}
\usepackage{booktabs}
\usepackage{geometry}
\usepackage{siunitx}
\usepackage{hyperref}
\usepackage{xcolor}
\usepackage{tcolorbox}
\usepackage{longtable}
\usepackage{enumitem}
\geometry{margin=2.5cm}

\title{Reported UAP Flight Profiles Matched Against\\
DFD $\psi$-Bubble Constraints:\\
\large A Quantitative Technology-Level Assessment}
\author{Agent 10 --- Iteration: $\psi$-Bubble Earth Operations}
\date{28 March 2026}

\begin{document}
\maketitle

\begin{abstract}
We perform a systematic, quantitative comparison between nine commonly reported
Unidentified Aerial Phenomena (UAP) flight characteristics and the operational
constraints of a $\psi$-bubble device as derived from Density Field Dynamics (DFD).
For each reported behaviour we compute the minimum magnetic field $B_{\min}$
required (with and without plasma-sheath density reduction), the stored energy,
the continuous power demand, and the implied technology level on a civilisation
scale.  The analysis draws directly on Agents~1--6 of this iteration and uses
only the established DFD equations: $\eta_c = \alpha/4$, the threshold field
$B_{\rm thr} = \sqrt{2\mu_0 \eta_c \rho c^2}$, Mode~I atmospheric thrust,
and the freefall geodesic (zero internal acceleration) property of the bubble.

\medskip\noindent
\textbf{Executive result:}
Five of the nine reported UAP behaviours (right-angle turns, no sonic boom,
no exhaust, luminous glow, EM interference) are \emph{natural, zero-cost
consequences} of any above-threshold $\psi$-bubble.  Two behaviours
(hypersonic acceleration, all-altitude operation) require $B \sim 20$--$100$\,T
with plasma-sheath assistance, placing them at or just beyond current
superconducting magnet technology.  Sea-level hover demands continuous
GW-class power.  Trans-medium (water-to-air) operation requires either
$B > 600{,}000$\,T without density reduction (impossible) or a self-sustaining
vapour/plasma envelope reducing water density by $\sim 10^5$, which is
physically conceivable but far beyond any known engineering.
\end{abstract}

\tableofcontents
\newpage

%=============================================================================
\section{Governing DFD Equations}\label{sec:governing}
%=============================================================================

All calculations use the following framework from the DFD unified theory
and Agents~1--6 of this iteration.

\subsection{Threshold Condition}

The $\psi$-bubble activates when:
\begin{equation}\label{eq:eta}
\eta \;\equiv\; \frac{U_{\rm EM}}{\rho\, c^2}
\;=\; \frac{B^2}{2\mu_0\,\rho\,c^2}
\;\geq\; \eta_c \;=\; \frac{\alpha}{4}
\;\approx\; 1.824 \times 10^{-3}.
\end{equation}

\subsection{Threshold Field at Given Density}

Solving for $B$:
\begin{equation}\label{eq:Bthr}
B_{\rm thr}(\rho)
= \sqrt{2\mu_0\,\eta_c\,\rho\,c^2}
= \sqrt{2 \times 1.257\times 10^{-6} \times 1.824\times 10^{-3}
    \times \rho \times 8.988\times 10^{16}}
\end{equation}
\begin{equation}
\boxed{B_{\rm thr}(\rho) = 20{,}329 \;\sqrt{\rho} \;\;\text{[T, with $\rho$ in kg/m$^3$]}}
\end{equation}

This gives:
\begin{center}
\begin{tabular}{lcc}
\toprule
\textbf{Environment} & $\boldsymbol{\rho}$ \textbf{(kg/m$^3$)} & $\boldsymbol{B_{\rm thr}}$ \textbf{(T)} \\
\midrule
Sea-level air        & 1.225            & 22{,}490 \\
10 km (cruising)     & 0.414            & 13{,}070 \\
30 km (stratosphere) & $1.84\times 10^{-2}$ & 2{,}757 \\
50 km (mesosphere)   & $1.03\times 10^{-3}$ & 652 \\
80 km (mesopause)    & $1.85\times 10^{-5}$ & 87.3 \\
100 km (K\'{a}rm\'{a}n line) & $5.6\times 10^{-7}$ & 15.2 \\
200 km (LEO)         & $2.5\times 10^{-10}$ & 0.32 \\
Space (interplanetary) & $\sim 10^{-20}$ & $\sim 10^{-7}$ \\
\midrule
Freshwater           & 1{,}000          & 642{,}800 \\
Seawater             & 1{,}025          & 651{,}800 \\
Plasma cavity ($\rho_{\rm local} \sim 10^{-5}$) & $10^{-5}$ & 2.03 \\
Steam envelope ($\rho \sim 0.6$) & 0.6  & 15{,}740 \\
\bottomrule
\end{tabular}
\end{center}

\subsection{Plasma-Sheath Density Reduction}

From Agent~1 and Agent~2: a 20\,T superconducting magnet generates magnetic
pressure $P_{\rm mag} = B^2/(2\mu_0) = 1.59\times 10^8$\,Pa, which is $\sim 1{,}570$
times sea-level atmospheric pressure.  This expels air from a cavity extending
to $r_{\rm cav} \approx 3.4\,R_s$ (Agent~2, Eq.~8), leaving a residual density
of order $\rho_{\rm res} \sim 10^{-5}$--$10^{-3}$\,kg/m$^3$ depending on the
completeness of expulsion and the rate of re-entry.

With plasma-sheath assistance (the expelled air ionises and is further
confined by the $\psi$-gradient), the effective density inside the bubble
can reach $\rho_{\rm eff} \sim 10^{-5}$\,kg/m$^3$ at sea level,
reducing $B_{\rm thr}$ from 22{,}490\,T to:
\begin{equation}
B_{\rm thr}^{\rm (plasma)} = 20{,}329\,\sqrt{10^{-5}} \approx 64\;\text{T.}
\end{equation}

Agent~2 demonstrated that the combined magnetic expulsion + plasma feedback
loop can sustain a low-density cavity with surface fields of
$B \sim 20$--$50$\,T, though with a continuous power demand of order
$\sim$\,TW\,m$^{-3}$ at sea level (bremsstrahlung replacement) or
$< 1$\,MW at altitudes $h \gtrsim 85$\,km.


%=============================================================================
\section{UAP Behaviour 1: Sea-Level Hover}\label{sec:hover}
%=============================================================================

\subsection{Reported Observation}

Stationary hovering over water or land at altitudes of $0$--$100$\,m,
lasting seconds to minutes, with no visible means of lift.

\subsection{DFD Requirements}

\paragraph{Threshold crossing.}
At sea level ($\rho = 1.225$\,kg/m$^3$), the naive threshold is
$B_{\rm thr} = 22{,}490$\,T.  With a plasma-sheath cavity reducing
$\rho_{\rm local}$ to $\sim 10^{-5}$\,kg/m$^3$: $B_{\rm thr} \approx 64$\,T.

With the optimistic bootstrap scenario from Agent~1 (self-sustaining
plasma feedback at the magnetic pressure boundary): $B \sim 20$--$50$\,T
may suffice.

\paragraph{Lift mechanism.}
Once above threshold, the device is in geodesic freefall within a modified
$n_{\rm eff}$ region.  To hover requires either:
\begin{itemize}
\item \textbf{Mode~I (ambient air ingestion):} Asymmetric $\psi$-gradient
  deflects infalling air downward.  Agent~4 computed $F_{\rm Mode\,I} \sim 5$\,MN
  at sea level for a 5\,m-radius shell with $B = 20$\,T.  This vastly exceeds
  the device weight ($\sim 10$\,kN for a 1{,}000\,kg craft).
\item \textbf{$n_{\rm eff}$ gradient:} The bubble itself modifies the local
  metric, creating an effective gravitational gradient.  The device rides this
  gradient.
\end{itemize}

\paragraph{Power requirement.}
\begin{itemize}
\item Maintaining the superconducting magnet: $\sim$\,kW (cryogenics).
\item Bremsstrahlung replacement for the plasma sheath at sea level:
  Agent~2 computed $P_{\rm brem} \sim 27$\,TW/m$^3$, giving
  $\sim 10^{13}$\,W for a $\sim 500$\,m$^3$ cavity volume.
\item \textbf{Implication:} Sea-level hover at steady state requires
  $\sim$\,GW to TW of continuous power for sheath maintenance.
\end{itemize}

\begin{tcolorbox}[colback=yellow!5!white, colframe=yellow!60!black,
  title=Sea-Level Hover Assessment]
\textbf{Without plasma sheath:} $B > 22{,}000$\,T.
\emph{Category (d): impossible with known physics.}\\[4pt]
\textbf{With plasma sheath:} $B \sim 20$--$64$\,T,
power $\sim$\,GW--TW continuous.\\
\emph{Category (c): advanced technology, far beyond current capability.}\\[4pt]
The magnetic field itself is within reach of present-day pulsed magnets
(45\,T sustained at NHMFL, 100\,T pulsed).  The bottleneck is the
\emph{continuous power} for sheath maintenance at sea-level density.
\end{tcolorbox}


%=============================================================================
\section{UAP Behaviour 2: Hypersonic Acceleration}\label{sec:hypersonic}
%=============================================================================

\subsection{Reported Observation}

Acceleration from stationary to Mach~5+ ($> 1{,}700$\,m/s) in $1$--$5$
seconds, implying $a \sim 350$--$1{,}700$\,m/s$^2$ ($\sim 35$--$170\,g$).

\subsection{DFD Analysis}

\paragraph{Acceleration mechanism.}
The fundamental DFD body-force acceleration is:
\begin{equation}
a = \frac{c^2}{2}\,\kappa\,\nabla\eta.
\end{equation}
With $\kappa = \eta_c = \alpha/4$ and $\nabla\eta$ set by the
gradient across the bubble of radius $r_b$:
\begin{equation}
a \sim \frac{c^2}{2} \times \frac{\alpha}{4} \times \frac{\eta}{r_b}
\sim \frac{(3\times 10^8)^2}{2} \times \frac{1.82\times 10^{-3}}{10}
\times \eta.
\end{equation}
For $\eta \sim \eta_c$ (just above threshold): $a \sim 10^{10}$\,m/s$^2$.

This is $\sim 10^9\,g$ --- far exceeding any reported UAP acceleration.
\textbf{Once above threshold, the acceleration available is essentially
unlimited on human scales.}

\paragraph{Energy budget.}
The kinetic energy to reach Mach~5 ($v = 1{,}700$\,m/s) for a 1{,}000\,kg device:
\begin{equation}
E_k = \tfrac{1}{2}mv^2 = \tfrac{1}{2}(1{,}000)(1{,}700)^2
= 1.45 \times 10^9\;\text{J} = 1.45\;\text{GJ}.
\end{equation}

The stored magnetic energy in a 50\,T, 10\,m$^3$ cavity is
$E_{\rm mag} = (B^2/2\mu_0)\times V = 9.95\times 10^9$\,J $\approx 10$\,GJ.
So a single acceleration to Mach~5 drains $\sim 15\%$ of stored energy
(if the energy comes from the field).  If the energy comes from ambient
air momentum transfer (Mode~I), no field energy is consumed.

\paragraph{Sustained hypersonic flight.}
At altitude ($h > 80$\,km), the threshold is easily crossed with $B > 87$\,T
(or much less with residual atmospheric density reduction).  The device can
accelerate essentially instantaneously and sustain hypersonic speeds
with negligible drag (the $\psi$-bubble deflects incident air around the
device; see Section~\ref{sec:boom}).

\begin{tcolorbox}[colback=green!5!white, colframe=green!60!black,
  title=Hypersonic Acceleration Assessment]
\textbf{At altitude ($> 80$\,km):} $B \sim 20$--$100$\,T sufficient.\\
\emph{Category (a)--(b): present to near-future technology.}\\[4pt]
\textbf{At sea level (with plasma sheath):} $B \sim 50$\,T + GW power.\\
\emph{Category (c): advanced technology.}\\[4pt]
The DFD framework provides \emph{more than enough} acceleration.
The reported 35--170$g$ values are trivial compared to the $\sim 10^9\,g$
theoretically available.  The limiting factor is not the physics but
the energy supply for maintaining the bubble.
\end{tcolorbox}


%=============================================================================
\section{UAP Behaviour 3: Instantaneous Right-Angle Turns}\label{sec:turns}
%=============================================================================

\subsection{Reported Observation}

Objects travelling at high speed execute $90^\circ$ turns with no
observable deceleration-acceleration phase and no structural deformation.

\subsection{DFD Analysis}

\paragraph{Freefall geodesic.}
This is the single most natural prediction of the $\psi$-bubble.
The device and its occupants are in \emph{freefall} within the bubble ---
they follow geodesics of the locally modified metric.  There is
\textbf{zero proper acceleration} experienced inside, regardless of
the coordinate acceleration seen by external observers.

\paragraph{Direction change mechanism.}
To change the direction of travel:
\begin{enumerate}
\item The asymmetry of the EM field (which sets $\nabla\eta$) is
  electronically reoriented --- e.g.\ by shifting currents in a
  segmented superconducting coil array.
\item The new $\nabla\eta$ defines a new preferred geodesic direction.
\item The device follows the new geodesic.
\item Switching time: limited by coil inductance and switching electronics.
  For a well-designed system: $\tau_{\rm switch} \sim 1$--$10$\,ms.
\end{enumerate}

\paragraph{No structural stress.}
Because the entire device (and its contents) follows the same geodesic,
there are \emph{no tidal forces} across the device (to first order in
$r_{\rm device}/r_{\rm bubble}$).  The structural integrity question is
identical to that of an astronaut in orbit: none, because everything is
in freefall.

\paragraph{No additional field or energy required.}
The right-angle turn requires no additional magnetic field strength
beyond what is already needed to maintain the bubble.  It requires only
that the EM asymmetry be electronically redirectable.

\begin{tcolorbox}[colback=green!5!white, colframe=green!60!black,
  title=Right-Angle Turn Assessment]
\textbf{Field requirement:} Same as for any bubble operation --- no extra.\\
\textbf{Technology:} Electronic coil switching ($\sim$\,ms timescale).\\
\emph{Category (a): this is a free, natural consequence of the bubble.}\\[4pt]
Right-angle turns at any speed are the \emph{signature prediction} of
a $\psi$-bubble device.  Any technology that creates a genuine modified-metric
bubble will exhibit this behaviour.  \textbf{Zero g-forces inside.}
\end{tcolorbox}


%=============================================================================
\section{UAP Behaviour 4: Trans-Medium Operation}\label{sec:transmedium}
%=============================================================================

\subsection{Reported Observation}

Objects enter the ocean from air (or emerge from water into air) with
no apparent change in speed or behaviour.  Some reports describe
extended underwater travel.

\subsection{DFD Analysis}

\paragraph{Threshold in water.}
For seawater ($\rho = 1{,}025$\,kg/m$^3$):
\begin{equation}
B_{\rm thr}^{\rm (water)} = 20{,}329\,\sqrt{1{,}025} = 651{,}000\;\text{T.}
\end{equation}
This is \textbf{utterly impossible} with any conceivable magnet technology.

\paragraph{With density reduction.}
However, the question is what the \emph{local} density inside the bubble is
when the device enters water.  Three scenarios:

\begin{enumerate}[label=(\roman*)]
\item \textbf{No preparation:} The device, already operating in air with a
  plasma sheath, plunges into water.  The water rushes in at the
  sheath boundary.  Local density jumps from $\sim 10^{-5}$ to
  $\sim 10^3$\,kg/m$^3$ in milliseconds.  The bubble collapses.
  \emph{Transit fails.}

\item \textbf{Steam/vapour envelope:} The intense EM field and plasma sheath
  flash-boil incoming water.  Steam density is $\rho_{\rm steam} \approx
  0.6$\,kg/m$^3$ at 100$^\circ$C and 1\,atm.  At higher temperatures,
  steam density drops further.  For superheated steam at $T \sim 10{,}000$\,K
  and pressure equilibrium with surrounding water ($\sim 1$\,atm at surface):
  \begin{equation}
  \rho_{\rm hot\,steam} \sim \frac{P\,\mu}{R\,T}
  = \frac{10^5 \times 0.018}{8.314 \times 10^4}
  \approx 0.022\;\text{kg/m$^3$.}
  \end{equation}
  Threshold: $B_{\rm thr} = 20{,}329 \times \sqrt{0.022} = 3{,}010$\,T.
  Still far too high.

  At $T \sim 10^6$\,K (fully ionised plasma):
  $\rho \sim 2\times 10^{-5}$\,kg/m$^3$.
  Threshold: $B_{\rm thr} = 20{,}329\times\sqrt{2\times 10^{-5}} = 91$\,T.
  This requires maintaining a $\sim 10^6$\,K plasma envelope against
  ocean water on all sides --- an extraordinary challenge.

\item \textbf{Supercavitation + plasma:} The device creates a
  supercavitating vapour cavity (like a torpedo) augmented by plasma
  heating.  The vapour pocket has $\rho \sim 10^{-2}$\,kg/m$^3$
  (hot steam at moderate temperatures), and the plasma-sheath bootstrap
  further reduces interior density.  If $\rho_{\rm eff} \sim 10^{-5}$:
  $B_{\rm thr} = 64$\,T.

  Power for continuous water vaporisation:
  To maintain a cavity of radius $R = 17$\,m in water requires
  vaporising the inrushing water at rate:
  \begin{equation}
  \dot{m} = \rho_w \times A \times v_{\rm collapse}
  \end{equation}
  where $v_{\rm collapse} \sim \sqrt{2P/\rho_w} \approx 14$\,m/s
  (Rayleigh collapse speed at 1\,atm).  For $A = 4\pi(17)^2 \approx
  3{,}600$\,m$^2$:
  $\dot{m} \approx 1{,}025 \times 3{,}600 \times 14 = 5.2\times 10^7$\,kg/s.

  Latent heat of vaporisation: $L = 2.26\times 10^6$\,J/kg.  Power:
  \begin{equation}
  P_{\rm vap} = \dot{m}\times L = 5.2\times 10^7 \times 2.26\times 10^6
  = 1.2\times 10^{14}\;\text{W} = 120\;\text{TW.}
  \end{equation}

  This is approximately \textbf{7 times total human energy production}
  ($\sim 18$\,TW).  Even at smaller radii the power is enormous.
\end{enumerate}

\begin{tcolorbox}[colback=red!5!white, colframe=red!60!black,
  title=Trans-Medium Assessment]
\textbf{Without density reduction:} $B > 650{,}000$\,T.\\
\emph{Category (d): impossible.}\\[4pt]
\textbf{With steam/plasma envelope:} $B \sim 64$--$100$\,T,
power $\sim 100$\,TW continuous.\\
\emph{Category (d): physically conceivable but power requirements are
$\sim 10\times$ human civilisation output.}\\[4pt]
Trans-medium operation is the \emph{hardest} UAP behaviour to reproduce
within DFD.  It requires either a field strength far beyond any known
physics, or a power source exceeding anything in human experience.
A civilisation capable of this would need to command $\sim 10^{14}$\,W
from a device of $\sim 10$\,m scale --- implying an energy density
source far beyond nuclear fission or fusion.  Antimatter annihilation
($\sim 10^{16}$\,J/kg) could in principle provide this, but the
engineering challenges of storing and releasing it are profound.
\end{tcolorbox}


%=============================================================================
\section{UAP Behaviour 5: All-Altitude Operation}\label{sec:altitude}
%=============================================================================

\subsection{Reported Observation}

Objects observed transitioning seamlessly from sea level to near-space
altitudes ($> 80$\,km), or appearing at arbitrary altitudes.

\subsection{DFD Analysis}

Agent~6 computed the minimum operating altitude as a function of $B$
(without plasma sheath):

\begin{center}
\begin{tabular}{cc}
\toprule
$B$ (T) & $h_{\min}$ (km) \\
\midrule
10  & 131.2 \\
20  & 119.4 \\
50  & 103.8 \\
100 & 92.1  \\
200 & 80.3  \\
500 & 64.7  \\
1{,}000 & 52.9 \\
\bottomrule
\end{tabular}
\end{center}

With a plasma-sheath cavity ($\rho_{\rm local} \sim 10^{-5}$\,kg/m$^3$),
the threshold is crossed at $B \sim 64$\,T at \emph{any} altitude
(including sea level), provided the sheath is maintained.

\paragraph{Altitude transition profile.}
The optimal operating profile is:
\begin{enumerate}
\item \textbf{Launch at $h > 80$\,km:}  With $B = 100$\,T, the threshold
  is crossed in the ambient atmosphere.  No plasma sheath needed.
  Power: $\sim$\,kW (cryogenics only).
\item \textbf{Descend with increasing field:}  As the device descends,
  atmospheric density rises exponentially.  The plasma sheath activates
  automatically (magnetic pressure exceeds atmospheric pressure for
  $B > 0.5$\,T at sea level).
\item \textbf{Sea-level operation:}  Full plasma-sheath mode.
  Power: $\sim$\,GW--TW.
\end{enumerate}

\begin{tcolorbox}[colback=green!5!white, colframe=green!60!black,
  title=All-Altitude Assessment]
\textbf{High altitude ($> 80$\,km):} $B \sim 20$--$100$\,T, power $\sim$\,kW.\\
\emph{Category (a)--(b): present to near-future technology.}\\[4pt]
\textbf{Mid-altitude (30--80\,km):} $B \sim 50$--$200$\,T with partial
plasma sheath, power $\sim$\,MW.\\
\emph{Category (b)--(c): near-future to advanced.}\\[4pt]
\textbf{Sea level:} $B \sim 20$--$64$\,T with full plasma sheath,
power $\sim$\,GW--TW.\\
\emph{Category (c)--(d): advanced to extreme.}
\end{tcolorbox}


%=============================================================================
\section{UAP Behaviour 6: No Sonic Boom}\label{sec:boom}
%=============================================================================

\subsection{Reported Observation}

Objects travelling at supersonic or hypersonic speeds produce no audible
sonic boom.

\subsection{DFD Analysis}

\paragraph{Standard sonic boom mechanism.}
A conventional supersonic object pushes air molecules ahead of it.  When
the object exceeds the local speed of sound, the displaced air cannot
move aside fast enough and piles up into a shock wave.

\paragraph{$\psi$-bubble mechanism.}
A $\psi$-bubble device does \emph{not push air}.  Instead:
\begin{enumerate}
\item The bubble creates a modified-metric region around the device.
\item Air molecules entering the bubble are deflected around the device
  by the $\psi$-gradient force, not by mechanical contact.
\item The deflection occurs smoothly over the bubble radius
  ($r_b \sim 10$--$20$\,m), not at a sharp surface.
\item From the air's perspective, it experiences a gentle gravitational
  lens effect --- it curves around the device and resumes its
  original trajectory behind.
\item There is no compression wave because there is no mechanical
  displacement of air ahead of the device.
\end{enumerate}

\paragraph{Condition for no boom.}
The bubble must be large enough that the deflection angle per unit
length is gentle:
\begin{equation}
\theta_{\rm deflect} \sim \frac{v_{\rm device}}{c} \times \frac{r_{\rm device}}{r_{\rm bubble}}.
\end{equation}
For $v = 5{,}000$\,m/s (Mach~15), $r_{\rm device} = 5$\,m,
$r_{\rm bubble} = 17$\,m:
$\theta \sim (5{,}000/3\times 10^8)(5/17) \sim 5\times 10^{-6}$\,rad.
This is negligible --- no shock forms.

\begin{tcolorbox}[colback=green!5!white, colframe=green!60!black,
  title=No Sonic Boom Assessment]
\textbf{Field requirement:} None beyond what is needed for the bubble itself.\\
\emph{Category (a): this is a free, automatic consequence of the bubble.}\\[4pt]
The absence of a sonic boom is a \emph{necessary prediction} of any
$\psi$-bubble device.  If a genuine modified-metric bubble is operating,
sonic booms are \emph{impossible} regardless of speed.
\end{tcolorbox}


%=============================================================================
\section{UAP Behaviour 7: No Visible Exhaust}\label{sec:exhaust}
%=============================================================================

\subsection{Reported Observation}

No rocket plume, jet exhaust, contrail, or any visible propellant
emission.

\subsection{DFD Analysis}

The $\psi$-bubble does not expel propellant.  Thrust comes from:
\begin{itemize}
\item \textbf{Mode~I:} Ambient air is gravitationally deflected
  rearward.  The ``exhaust'' is ordinary air at near-ambient temperature,
  not a hot combustion plume.  It would be invisible.
\item \textbf{Mode~II:} At high altitude or in vacuum, thrust comes
  from the metric gradient itself.  No mass is expelled.
\item \textbf{Plasma sheath:} The plasma sheath is \emph{confined}, not
  expelled.  It may glow (Section~\ref{sec:glow}) but does not form
  an exhaust trail.
\end{itemize}

\begin{tcolorbox}[colback=green!5!white, colframe=green!60!black,
  title=No Exhaust Assessment]
\textbf{Field requirement:} None beyond the bubble itself.\\
\emph{Category (a): free consequence of bubble propulsion.}\\[4pt]
The absence of exhaust is \emph{definitional} for a $\psi$-bubble device.
Any genuine DFD propulsion system produces no visible exhaust by design.
\end{tcolorbox}


%=============================================================================
\section{UAP Behaviour 8: Luminous Glow (Variable Colour)}\label{sec:glow}
%=============================================================================

\subsection{Reported Observation}

Objects emit light --- often described as white, blue-white, orange,
or red --- sometimes shifting colour.  The glow appears to envelop the
entire object.

\subsection{DFD Analysis}

\paragraph{Plasma sheath emission.}
A sea-level or low-altitude $\psi$-bubble maintains a plasma sheath at
temperatures of $10^4$--$10^6$\,K (Agent~2).  This plasma radiates:

\begin{itemize}
\item \textbf{$T \sim 5{,}000$--$6{,}000$\,K:} White/yellow-white
  (solar-type blackbody spectrum).
\item \textbf{$T \sim 10{,}000$\,K:} Blue-white
  (recombination-dominated nitrogen/oxygen plasma).
\item \textbf{$T \sim 30{,}000$--$100{,}000$\,K:} Deep blue/UV
  (bremsstrahlung-dominated).
\item \textbf{$T \sim 3{,}000$\,K:} Orange-red
  (cooler outer sheath boundary).
\end{itemize}

\paragraph{Colour variation.}
The plasma temperature depends on the operating power and the
ambient density.  A device adjusting its field strength, altitude, or
power level would naturally produce colour shifts:
\begin{itemize}
\item Increasing power $\rightarrow$ hotter plasma $\rightarrow$ bluer.
\item Decreasing power / higher altitude $\rightarrow$ cooler / thinner
  plasma $\rightarrow$ redder, dimmer, eventually invisible.
\item Transition from sea level to altitude: glow fades as
  atmospheric density (and hence plasma density) drops.
\end{itemize}

\paragraph{High-altitude glow.}
At $h > 100$\,km, residual atmospheric molecules excited by the EM
field emit auroral-type emissions (green from oxygen at 557.7\,nm,
red at 630\,nm).

\begin{tcolorbox}[colback=green!5!white, colframe=green!60!black,
  title=Luminous Glow Assessment]
\textbf{Field requirement:} None beyond the bubble + plasma sheath.\\
\emph{Category (a): natural byproduct of plasma sheath operation.}\\[4pt]
A luminous glow of variable colour is a \emph{firm prediction} of any
atmospheric $\psi$-bubble device.  The colour encodes the plasma
temperature and hence the operating state.  The reported colour
variations (white, blue, orange, red) are exactly what DFD predicts
for different power levels and altitudes.
\end{tcolorbox}


%=============================================================================
\section{UAP Behaviour 9: EM Interference}\label{sec:emi}
%=============================================================================

\subsection{Reported Observation}

Electronics malfunction in the vicinity of UAP: radar dropout,
compass spin, radio static, vehicle electrical failures.

\subsection{DFD Analysis}

\paragraph{Magnetic field environment.}
A $\psi$-bubble device with $B_0 = 50$\,T at the shell surface
has a dipole field that falls off as $(R/r)^3$.  At distance $d$
from the shell:
\begin{equation}
B(d) = B_0\left(\frac{R}{R+d}\right)^3.
\end{equation}

At $d = 100$\,m from a 5\,m-radius device:
$B(100) = 50\times(5/105)^3 = 50\times 1.08\times 10^{-4} = 5.4\;\text{mT}$.

At $d = 500$\,m: $B(500) = 50\times(5/505)^3 = 50\times 9.7\times 10^{-7}
= 48\;\mu$T $\approx$ Earth's field.

\paragraph{Effects on electronics.}
\begin{itemize}
\item At $< 50$\,m: $B > 50$\,mT.  Strong enough to saturate magnetic
  sensors, disrupt compass operation, induce currents in wiring,
  corrupt magnetic storage media.
\item At $50$--$200$\,m: $B \sim 0.05$--$5$\,mT.  Can affect
  sensitive instruments, cause radar anomalies, induce noise in
  radio receivers.
\item At $> 500$\,m: field comparable to Earth's.  Minimal direct
  magnetic effect, but the \emph{time-varying} field from a moving
  device induces $\mathcal{E} = -d\Phi/dt$ in any conducting loop.
  For a device passing at 1{,}000\,m/s at 200\,m distance, the
  rate of field change can produce significant induced voltages.
\end{itemize}

\paragraph{Plasma EM emission.}
The plasma sheath also radiates broadband EM noise (synchrotron,
bremsstrahlung, cyclotron emission), which would appear as
wideband static across radio frequencies.

\paragraph{$\psi$-gradient effects on photon propagation.}
Within the modified-metric region, the local speed of light is
altered: $c_{\rm local} = c/n_{\rm eff}$.  Radar pulses entering
the bubble region experience refraction and possible total internal
reflection, explaining radar dropout and anomalous returns.

\begin{tcolorbox}[colback=green!5!white, colframe=green!60!black,
  title=EM Interference Assessment]
\textbf{Field requirement:} None beyond the bubble itself.\\
\emph{Category (a): unavoidable consequence of the fields involved.}\\[4pt]
EM interference is an \emph{unavoidable side effect} of any device
operating with $B > 20$\,T fields.  The combination of static dipole
field, time-varying induction from motion, broadband plasma emission,
and modified photon propagation within the bubble produces exactly
the reported effects: compass disruption, radar anomalies, radio
static, and electrical failures.
\end{tcolorbox}


%=============================================================================
\section{Master Comparison Table}\label{sec:master}
%=============================================================================

\begin{longtable}{p{2.8cm}cccp{3.5cm}c}
\toprule
\textbf{UAP Behaviour} &
\makecell{$B_{\min}$\\(no sheath)} &
\makecell{$B_{\min}$\\(w/ sheath)} &
\textbf{Power} &
\textbf{Limiting Factor} &
\textbf{Category} \\
\midrule
\endfirsthead
\toprule
\textbf{UAP Behaviour} &
\makecell{$B_{\min}$\\(no sheath)} &
\makecell{$B_{\min}$\\(w/ sheath)} &
\textbf{Power} &
\textbf{Limiting Factor} &
\textbf{Category} \\
\midrule
\endhead
Sea-level hover     & 22,500\,T & 20--64\,T & GW--TW  & Power for sheath  & (c)--(d) \\
Hypersonic accel.   & 87\,T$^*$  & 20--50\,T & MW--GW  & Energy budget     & (a)--(c)$^\dagger$ \\
Right-angle turns   & ---        & ---       & $\sim 0$ & Coil switching    & (a) \\
Trans-medium        & 652,000\,T & 64--100\,T & 100+ TW & Power (water vap.) & (d) \\
All altitudes       & 15--22,500\,T & 20--64\,T & kW--TW & Varies w/ altitude & (a)--(d)$^\dagger$ \\
No sonic boom       & ---        & ---       & $\sim 0$ & None (automatic)  & (a) \\
No exhaust          & ---        & ---       & $\sim 0$ & None (by design)  & (a) \\
Luminous glow       & ---        & ---       & $\sim 0$ & None (automatic)  & (a) \\
EM interference     & ---        & ---       & $\sim 0$ & None (automatic)  & (a) \\
\bottomrule
\multicolumn{6}{l}{\footnotesize $^*$At 80\,km altitude. $^\dagger$Depends on altitude.}\\
\multicolumn{6}{l}{\footnotesize Categories: (a) present tech 1--50\,T;
  (b) near-future 50--200\,T; (c) advanced 200--10,000\,T; (d) impossible/$>$10,000\,T.}
\end{longtable}


%=============================================================================
\section{Technology-Level Classification}\label{sec:civlevel}
%=============================================================================

\subsection{Magnetic Field Technology Scale}

\begin{center}
\begin{tabular}{cll}
\toprule
\textbf{$B$ (T)} & \textbf{Technology} & \textbf{Status} \\
\midrule
1     & Permanent magnets, simple solenoids    & Available \\
10    & Superconducting MRI magnets             & Available \\
20    & Research SC magnets (LHC dipoles)        & Available \\
45    & DC hybrid magnets (NHMFL record)         & Available (single facility) \\
50    & Reinforced pulsed (non-destructive)       & Near-term \\
100   & Pulsed magnets (destructive, $\sim$ms)    & Available (single-shot) \\
200   & Reinforced multi-shot pulsed              & Near-future \\
500   & Explosive flux compression                & Available (destructive) \\
1{,}000  & Advanced flux compression             & Developmental \\
2{,}800  & Record: Z-machine (Sandia, 2012)      & Single-shot, ns pulse \\
10{,}000 & Hypothetical nuclear-magnetic          & Not demonstrated \\
100{,}000+ & Beyond known engineering             & Speculative \\
\bottomrule
\end{tabular}
\end{center}

\subsection{Power Source Technology Scale}

\begin{center}
\begin{tabular}{cll}
\toprule
\textbf{Power} & \textbf{Source} & \textbf{Status} \\
\midrule
kW    & Batteries, fuel cells                    & Available \\
MW    & Small nuclear reactor (SMR)              & Available \\
GW    & Large nuclear reactor / fusion (ITER)    & Near-future \\
TW    & Total human civilisation output          & Aggregate only \\
100\,TW & $\sim 6\times$ human civilisation      & Far beyond current \\
\bottomrule
\end{tabular}
\end{center}

\subsection{Where Each UAP Behaviour Falls}

\begin{enumerate}[label=\textbf{(\alph*)}]
\item \textbf{Present technology (1--50\,T, kW--MW power):}
  \begin{itemize}
  \item Right-angle turns (free with any bubble)
  \item No sonic boom (free)
  \item No exhaust (free)
  \item Luminous glow (free)
  \item EM interference (free)
  \item High-altitude ($>80$\,km) hover and acceleration (20--100\,T)
  \end{itemize}

\item \textbf{Near-future technology (50--200\,T, MW--GW power):}
  \begin{itemize}
  \item Mid-altitude (30--80\,km) operation with partial plasma sheath
  \item Sustained hypersonic flight at moderate altitudes
  \end{itemize}

\item \textbf{Advanced technology requiring new engineering (50--100\,T sustained,
  GW--TW power):}
  \begin{itemize}
  \item Sea-level hover (field feasible, power supply is the bottleneck)
  \item Sea-level hypersonic acceleration
  \item All-altitude seamless transition
  \end{itemize}

\item \textbf{Beyond current framework ($>$10,000\,T or $>$100\,TW):}
  \begin{itemize}
  \item Trans-medium (water-to-air) without density reduction envelope
  \item Sustained underwater operation at any speed
  \end{itemize}
\end{enumerate}


%=============================================================================
\section{Honest Assessment and Conclusions}\label{sec:conclusions}
%=============================================================================

\subsection{What DFD Explains Naturally (No Strain)}

Five of nine reported UAP behaviours --- right-angle turns, no sonic boom,
no exhaust, luminous glow, and EM interference --- are \emph{zero-cost
automatic consequences} of any above-threshold $\psi$-bubble.  They require
no additional technology beyond the bubble itself.  These five are, in fact,
firm \emph{predictions}: if DFD is correct and someone builds a $\psi$-bubble,
these behaviours \emph{must} occur.

\subsection{What DFD Explains With Strain (Advanced Technology)}

Two behaviours --- hypersonic acceleration and all-altitude operation ---
are fully consistent with DFD but require sustained magnetic fields of
20--100\,T and continuous power sources ranging from kW (high altitude)
to TW (sea level).  The fields are within or near current technology;
the power sources are the bottleneck, requiring at minimum a compact
GW-class power source (comparable to a large nuclear reactor) compressed
into a $\sim 10$\,m device.

Sea-level hover adds the further constraint of continuous plasma-sheath
maintenance against sea-level air density, demanding TW-class power.

\subsection{What DFD Cannot Easily Explain}

Trans-medium operation (water-to-air) is the single most challenging
reported behaviour.  Without a self-sustaining density reduction mechanism
in water, the required field strength ($>$650,000\,T) is beyond any
conceivable technology.  With a plasma/steam envelope, the field requirement
drops to 64--100\,T but the power requirement (100+\,TW for continuous
water vaporisation) exceeds total human civilisation output by an order
of magnitude.

\subsection{Implications for UAP Origin}

\begin{tcolorbox}[colback=blue!5!white, colframe=blue!60!black,
  title=Summary Verdict]
\begin{itemize}
\item \textbf{If UAP reports are accurate}, the DFD framework
  provides a \emph{consistent physical mechanism} for 7 of 9 reported
  behaviours with technology levels ranging from present-day to
  advanced-but-not-impossible.

\item \textbf{Trans-medium operation} requires either (a) physics
  beyond the current DFD framework, (b) a density reduction mechanism
  not yet analysed, or (c) a power source of extraordinary energy
  density ($\sim 10^{14}$\,W from a 10\,m device, implying
  $\sim 10^{10}$\,W/m$^3$).  Antimatter annihilation
  ($\sim 9\times 10^{16}$\,J/kg) could in principle supply this.

\item \textbf{The technology level implied by the full UAP capability
  set} (including trans-medium) corresponds to a civilisation that has
  mastered:
  \begin{enumerate}
  \item Sustained 50--100\,T superconducting magnets (near-future for us).
  \item Compact GW--TW power sources ($\sim$\,centuries ahead of us).
  \item Water-density reduction via plasma envelopes ($\sim$\,unknown
    timeline).
  \end{enumerate}

\item \textbf{On the Kardashev scale:} The full UAP capability set
  (especially trans-medium) implies a civilisation well above Type~0
  (humanity $\approx 0.73$) but not yet Type~I.  Roughly Type~0.9--1.0,
  commanding $\sim 10^{14}$--$10^{16}$\,W of directed power.

\item \textbf{DFD does not require ``new physics'' for UAP.}
  The framework is derivable from standard field theory with the
  single addition of the $\psi$-field vacuum coupling.  What is
  required is \emph{advanced engineering}: higher fields, larger
  power sources, better plasma confinement --- all extrapolations
  of known physics, not violations of it.
\end{itemize}
\end{tcolorbox}

\subsection{Falsifiability}

This analysis makes several falsifiable predictions:
\begin{enumerate}
\item Any genuine $\psi$-bubble device \emph{must} produce EM interference
  at ranges $< 500$\,m (dipole field $\geq$ Earth's field).
\item The luminous glow \emph{must} be thermal/plasma in character
  (not coherent laser-like emission), with a blackbody-like spectrum.
\item Right-angle turns \emph{must} be accompanied by no shock wave
  and no debris field.
\item A device operating at sea level \emph{must} displace nearby air
  (hurricane-force winds at the bubble boundary from magnetic pressure
  expulsion).
\item Trans-medium transition \emph{must} produce a massive steam
  explosion or cavitation event at the water entry/exit point.
\end{enumerate}

If careful UAP observations violate any of these, the DFD $\psi$-bubble
model would be ruled out as the mechanism.

\end{document}
